%MB9T5
%-------------------------------------------------
\begin{vidu}
Một sợi dây đàn hồi có hai đầu cố định $A,\ B$ cách nhau $L = 72\ cm$. Trên dây đã hình thành sóng dừng, biên độ tại các bụng là $A_b = 4\ cm$. $O$ là nút gần nhất kế bên một bụng $B$. Lấy $O$ làm gốc tọa độ và chiều dương từ $O$ ra xa $B$.
%----------- Ý 1 ----------
\item[1.] Điểm $M$ có biên độ $A_M = 0,5A_b$ và cách nút $O$ một đoạn $d = 2\ cm$.  
Bước sóng của sóng trên dây là
\choice
{18 cm}
{36 cm}
{\True 24 cm}
{30 cm}	
\item[2.] Hai điểm $P,\;Q$ cùng ở nửa dương của trục $x$ với toạ độ $x_P = 6\ cm$, $x_Q = 9\ cm$. Tỉ số biên độ $\dfrac{A_P}{A_Q}$ xấp xỉ
\choice
{$\dfrac12$}
{$\sqrt{3}$}
{\True $\sqrt{2}$}
{$\dfrac{\sqrt{3}}{2}$}
\item[3.] Khoảng cách ngắn nhất $MN$ giữa hai điểm gần nhau nhất có biên độ lần lượt  
$A_M = 2,6\ cm$ và $A_N = 2,1\ cm$ là
\choice
{$\dfrac{\lambda}{16}$}
{$\dfrac{\lambda}{12}$}
{\True $\dfrac{\lambda}{20}$}
{$\dfrac{\lambda}{8}$}
%----------- Ý 4 ----------
\item[4.] Gọi $C$ là trung điểm của đoạn nối một nút và bụng lân cận. Biết $CB = 3\ cm$. Thời gian ngắn nhất giữa hai lần $C$ và $B$ có cùng li độ là $\Delta t_{\min}=0,10\ s$.  
Tốc độ truyền sóng trên dây bằng
\choice
{0,90 m/s}
{1,80 m/s}
{2,40 m/s}
{\True 1,20 m/s}
%----------- Ý 5 ----------
\item[5.] Trên toàn bộ chiều dài dây $AB$, có bao nhiêu điểm dao động với biên độ $1,2\ cm$?
\choice
{6}
{8}
{10}
{\True 12}
%----------- Ý 6 ----------
\item[6.] Hai điểm $M_1,\,M_2$ cách nút $O$ những đoạn $x_1 = 4\ cm$ và $x_2 = 6\ cm$ (cùng một phía).  
Tại thời điểm dây đi qua vị trí cân bằng (li độ $u=0$), tỉ số vận tốc dao động $\dfrac{v_1}{v_2}$ là
\choice
{$\dfrac{\sqrt{6}}{3}$}
{$\sqrt{3}$}
{\True $\dfrac{\sqrt{3}}{2}$}
{$\dfrac{1}{\sqrt{2}}$}
\loigiai{
\begin{itemize}
%----- Giải 1 -----
\item[\textbf{1.}] $A_M = A_b\bigl|\sin\dfrac{2\pi d}{\lambda}\bigr|=0,5A_b  
\;\Rightarrow\; \sin\dfrac{2\pi d}{\lambda}=0,5  
\;\Rightarrow\; \dfrac{2\pi d}{\lambda}=\dfrac{\pi}{6}\;(\text{góc nhỏ nhất})  
\;\Rightarrow\; \lambda=12d = 24\ \text{cm}.$
%----- Giải 2 -----
\item[\textbf{2.}] $A(x)=A_b\bigl|\sin\dfrac{2\pi x}{\lambda}\bigr|,\;\lambda=24\ \text{cm}.  
\;A_P\propto\bigl|\sin\dfrac{2\pi\cdot6}{24}\bigr|=\bigl|\sin\dfrac{\pi}{2}\bigr|=1.  
\;A_Q\propto\bigl|\sin\dfrac{2\pi\cdot9}{24}\bigr|=\bigl|\sin\dfrac{3\pi}{4}\bigr|=\dfrac{\sqrt{2}}{2}.  
\;\dfrac{A_P}{A_Q}=\sqrt{2}.$
%----- Giải 3 -----
\item[\textbf{3.}] $A(x)=A_b\bigl|\cos\dfrac{2\pi x}{\lambda}\bigr|$ (lấy gốc tại bụng).  
Giải $\cos\dfrac{2\pi x_M}{\lambda}=\dfrac{2,6}{4}\;(=0,65),\  
\cos\dfrac{2\pi x_N}{\lambda}=\dfrac{2,1}{4}\;(=0,525)\;  
\Longrightarrow\; MN=\dfrac{\lambda}{20}.$
%----- Giải 4 -----
\item[\textbf{4.}] Với $CB=3\ \text{cm}\Rightarrow AB = 4CB = 12\ \text{cm}$.  
Thời gian cực tiểu C và B cùng li độ: $\Delta t_{\min}=\dfrac{T}{2}\Rightarrow T=0,20\ \text{s}$.  
$v=\dfrac{\lambda}{T}=\dfrac{24}{0,20}=120\ \text{cm/s}=1,20\ \text{m/s}.$
%----- Giải 5 -----
\item[\textbf{5.}] $A=1,2\ \text{cm}\Rightarrow \dfrac{A}{A_b}=0,3\  
\Longrightarrow \bigl|\sin\dfrac{2\pi x}{\lambda}\bigr|=0,3$.  
Trong mỗi nửa bước sóng $\lambda/2$ có 2 điểm thỏa mãn.  
Chiều dài dây: $L/\bigl(\lambda/2\bigr)=72/12=6$ nửa bước sóng  
$\Rightarrow$ số điểm $=6\times2=12$.  
%----- Giải 6 -----
\item[\textbf{6.}] Tại $u=0$, tốc độ tức thời $v=\omega A(x)$.  
$\displaystyle\dfrac{v_1}{v_2}=\dfrac{A(x_1)}{A(x_2)}  
=\dfrac{\sin\dfrac{2\pi\cdot4}{24}}{\sin\dfrac{2\pi\cdot6}{24}}  
=\dfrac{\sin\dfrac{\pi}{3}}{\sin\dfrac{\pi}{2}}=\dfrac{\sqrt{3}}{2}.$
\end{itemize}}
\end{vidu}



\loai{Biên độ cách nút 1 đoạn x}
\begin{vidu}%[3P2-9.2K][Loai1] 
\nguon{QG - 2018} Trên một sợi dây đàn hồi đang có sóng dừng với biên độ dao động của các điểm bụng là a. M là một phần tử dây dao động với biên độ 0,5a. Biết vị trí cân bằng của M cách điểm nút gần nó nhất một khoảng 2 cm. Sóng truyền trên dây có bước sóng là
\choice
{\True 24 cm}
{12 cm}
{16 cm}
{3 cm}
\loigiai{
\begin{itemize}
\item Biên độ của 1 phần tử M sóng dừng cách nút sóng đoạn d: $a_M=a\left|{\sin \dfrac{2\pi d}{\lambda}}\right|$.
\item Theo đề suy ra: $\dfrac{a_M}{a}=\left|\sin (\dfrac{2\pi d}{\lambda})\right|=\dfrac{0.5a}{a}=\dfrac{1}{2}\Rightarrow\dfrac{2\pi d}{\lambda}=\dfrac{\pi}{6}\Rightarrow\lambda =12d=12.2=24\ cm$.
\end{itemize}}
\end{vidu} \haidong{20}


\loai{Tỉ số biên độ dao động}
\begin{vidu}%[3P2-9.2K][Loai1] 
\nguon{QG - 2018} Một sợi dây đàn hồi căng ngang với đầu A cố định đang có sóng dừng. M và N là hai phần tử dao động điều hòa có vị trí cân bằng cách đầu A những đoạn lần lượt là 16 cm và 27 cm. Biết sóng truyền trên dây có bước sóng 24 cm. Tỉ số giữa biên độ dao động của M và biên độ dao động của N là
\choice
{$\dfrac{\sqrt{6}}{3}$}
{$\dfrac{\sqrt{3}}{2}$}
{$\dfrac{\sqrt{3}}{3}$}
{\True $\dfrac{\sqrt{6}}{2}$}
\loigiai{
\begin{itemize}
\item Đầu A cố định là một nút sóng nên ta có: $\dfrac{2a\left|{sin\left(\dfrac{2\pi.16}{24}\right)}\right|}{2a\left|{sin\left(\dfrac{2\pi.27}{24}\right)}\right|}=\dfrac{\sqrt{6}}{2}$.
\end{itemize}
}
\end{vidu} \haidong{20}

\loai{Tính khoảng cách ngắn nhất giữa 2 điểm M, N có biên độ}
\begin{vidu}%[3P2-9.2G][Loai1]
Trên dây AB có sóng dừng với bước sóng $\lambda$, biết bụng sóng có biên độ 4 cm tại vị trí M trên dây AB có biên độ $2\sqrt{3}$ cm; N là vị trí trên dây AB gần M nhất có biên độ $2\sqrt{2}$ cm. Khoảng cách MN bằng
\choice
{$\dfrac{\lambda }{12}$}
{$\dfrac{\lambda }{6}$}
{$\dfrac{5\lambda }{24}$}
{\True $\dfrac{\lambda }{24}$}
\loigiai{
\begin{itemize}
\item Chọn bụng gần M, N nhất làm gốc tọa độ khi đó (với $0 < x \le$  $\dfrac{\lambda }{4}$ ):\\ ${A_M}={{A}_{b}}. \left| \cos \dfrac{2\pi {{x}_{M}}}{\lambda } \right|\Rightarrow 2\sqrt{3}=4\left| \cos \dfrac{2\pi {{x}_{M}}}{\lambda } \right|\Rightarrow {{x}_{M}}=\dfrac{\lambda }{12}$.
\item Tương tự: ${A_N}={{A}_{b}}. \left| \cos \dfrac{2\pi {{x}_{N}}}{\lambda } \right|\Rightarrow 2\sqrt{2}=4\left| \cos \dfrac{2\pi {{x}_{N}}}{\lambda } \right|\Rightarrow {{x}_{N}}=\dfrac{\lambda }{8}$. 
\item Vậy $MN = \dfrac{\lambda }{8}-\dfrac{\lambda }{12}=\dfrac{\lambda }{24}$.
\end{itemize}
}
\end{vidu} \haidong{20}
\loai{Trạng thái giữa hai điểm trên cùng 1 bó sóng}
\begin{vidu}%[3P2-9.2K][Loai1]
Sóng dừng trên dây nằm ngang. Trong cùng bó sóng, A là nút, B là bụng, C là trung điểm AB. Biết $CB = 4$ cm. Thời gian ngắn nhất giữa hai lần C và B có cùng li độ là 0,13 s. Tốc độ truyền sóng trên dây là
\choice
{\True 1,23 m/s}
{2,46 m/s}
{3,24 m/s}
{0,98 m/s}
\loigiai{
\begin{itemize}
\item B có cùng li độ chỉ khi chúng ở vị trí cân bằng do đó khoảng thời gian ngắn nhất là:\\ ${{t}_{\min }}=\dfrac{T}{2}=0. 13\ s\Rightarrow T=0. 26\ s$.
\item Ta có $\lambda =4AB=8CB=8.4=32\ cm\Rightarrow v=\dfrac{\lambda }{T}=\dfrac{32}{0. 26}\approx 123\ cm/s$.
\end{itemize}
}
\end{vidu} \haidong{20}
\loai{Trạng thái của các điểm ở cùng thời điểm}
\begin{vidu}%[3P2-9.2G][Loai1]
Một sóng dừng trên dây có bước sóng $\lambda$ và N là một nút sóng. Hai điểm $M_1,\ M_2$ nằm về hai phía của N và có vị trí cân bằng cách N những đoạn lần lượt là $\dfrac{\lambda }{8}$ và $\dfrac{\lambda }{12}$. Ở cùng một thời điểm mà hai phần tử tại đó có li độ khác không thì tỉ số giữa li độ của $M_1$ so với $M_2$ là
\choice
{\True $\dfrac{u_1}{u_2}=-\sqrt{2}$}
{$\dfrac{u_1}{u_2}=\dfrac{1}{\sqrt{3}}$}
{$\dfrac{u_1}{u_2}=\sqrt{2}$}
{$\dfrac{u_1}{u_2}=-\dfrac{1}{\sqrt{3}}$}
\loigiai{
\begin{itemize}
\item Biểu thức của sóng dừng tại điểm M cách nút N: $NM = d$ . 
\item Chọn gốc tọa độ tại N. Ta có
$d_1 = NM_1 = -  $$\dfrac{\lambda }{8}$; $d_2 = NM_2 =  $$\dfrac{\lambda }{12}$. 
\item Biên độ của điểm $M_1$: $A_1 = 2a\left|\sin\left( \dfrac{2\pi \dfrac{\lambda}{8}}{\lambda }\right)\right|= a \sqrt{2}$  cm.
\item Biên độ của điểm $M_2$: $A_2 = 2a\left|\sin\left( \dfrac{2\pi \dfrac{\lambda}{12}}{\lambda }\right)\right|= a$  cm.
\item Do hai điểm ở trên hai bó sóng kề nhau nên dao động tại hai điểm ngược pha. 
\item Do đó tỉ số giữa li độ của $M_1$ so với $M_2$ là:  $\dfrac{u_1}{u_2} =  -\dfrac{A_1}{A_2}= - \sqrt{2}$.
\end{itemize}
}
\end{vidu} \haidong{20}
\loai{Tỉ số vận tốc}
\begin{vidu}%[3P2-9.2G][Loai1]
Một sóng dừng trên dây có bước sóng $\lambda$ và I là một nút sóng. Hai điểm $M_1$, $M_2$ nằm cùng một phía với I và có vị trí cân bằng cách I những đoạn lần lượt là $\dfrac{\lambda }{6}$ và $\dfrac{\lambda }{4}$. Khi dây không duỗi thẳng thì tỉ số giữa vận tốc của $M_1$ so với $M_2$ là
\choice
{$\dfrac{v_1}{v_2}=\dfrac{\sqrt{6}}{3}$}
{$\dfrac{v_1}{v_2}=-\dfrac{\sqrt{6}}{3}$}
{$\dfrac{v_1}{v_2}=\dfrac{\sqrt{6}}{2}$}
{\True $\dfrac{v_1}{v_2}=\dfrac{\sqrt{3}}{2}$}
\loigiai{
\begin{itemize}
\item Vì hai điểm nằm trên cùng một bó sóng nên dao động cùng pha. Do đó:\\ \begin{align*}\dfrac{v_1}{v_2}=\dfrac{A_1}{A_2}\dfrac{\sin \dfrac{2\pi x_1}{\lambda }}{\sin \dfrac{2\pi x_2}{\lambda }}=\dfrac{\sin \dfrac{\pi }{3}}{\sin \dfrac{\pi }{2}}=\dfrac{\sqrt{3}}{2}.\end{align*}
\end{itemize}
}
\end{vidu} \haidong{20}
\loai{Số điểm trên dây dao động với biên độ trung gian}
\begin{vidu}%[3P2-9.2K][Loai1]
\nguon{QG - 2018}
Một sợi dây đàn hồi dài 30 cm có hai đầu cố định. Trên dây đang có sóng dừng. Biết sóng truyền trên dây với bước sóng 20 cm và biên độ dao động của điểm bụng là 2 cm. Số điểm trên dây mà phần tử tại đó dao động với biên độ 6 mm là
\choice
{8}
{\True 6}
{3}
{4}
\loigiai{
\begin{itemize}
\item Ta có: $\dfrac{AB}{\lambda/2}=3\Rightarrow$ trên dây có 3 bó sóng, mỗi bó có 2 phần tử dao động với biên độ 6 mm.
\end{itemize}
}
\end{vidu} \haidong{20}

\loai{Khoảng cách lớn nhất - nhỏ nhất}
\begin{vidu}%[3P2-9.2G][Loai1]
Một sợi dây AB dài 24 cm, hai đầu cố định, đang có sóng dừng với hai bụng sóng. Khi dây duỗi thẳng, M và N là hai điểm trên dây chia sợi dây thành ba đoạn bằng nhau. Tỉ số khoảng cách lớn nhất và nhỏ nhất giữa hai điểm M và N trong quá trình sợi dây dao động là 1,25. Biên độ dao động bụng sóng là
\choice
{4 cm}
{5 cm}
{\True $2\sqrt{3}$ cm}
{$3\sqrt{3}$ cm}
\loigiai{
\begin{itemize}
\item Khoảng cách nhỏ nhất khi MN ở vị trí cân bằng là $\dfrac{AB}{3}=\dfrac{24}{3}=8\ cm$.
\item Gọi trung điểm MN là O (chính là 1 nút) khi đó $OM = 4$ cm.
\item Khoảng cách lớn nhất khi MN ở biên là: $8.1,25 = 10$ cm khi đó $OM = 10/2 = 5$ cm.
\item Vì M, N luôn dao động theo phương vuông góc phương truyền sóng, áp dụng định lý Pitago để tim biên độ dao động của M: $A_M=\sqrt{5^2-4^2}=3\ cm\Rightarrow A_b=\dfrac{A_M}{\left| \sin \dfrac{2\pi x_M}{\lambda} \right|}=\dfrac{3}{\left| \sin \dfrac{2\pi . 4}{24} \right|}=2\sqrt{3}\ cm$.
\end{itemize}
}
\end{vidu} \haidong{20}
\loai{Các điểm cùng biên độ cách đều nhau}
\begin{vidu}%[3P2-9.2K][Loai1]
Không xét các điểm bụng hoặc nút, quan sát thấy những điểm có cùng biên độ và ở gần nhau nhất thì đều cách đều nhau 15 cm. Bước sóng trên dây có giá trị bằng
\choice
{30 cm}
{\True 60 cm}
{90 cm}
{45 cm}
\loigiai{
\begin{itemize}
\item Các điểm có cùng biên độ, gần nhau nhất cách đều nhau thì cách nhau một khoảng $\dfrac{\lambda}{4} =15$ cm. Vậy  $\lambda =60$ cm.
\end{itemize}
}
\end{vidu} \haidong{20}
\loai{Khoảng thời gian ngắn nhất giữa hai lần li độ dao động của phần tử tại B bằng biên độ dao động của phần tử tại C}
\begin{vidu}%[3P2-9.2G][Loai1]
\nguon{MH - 2018} Một sợi dây đàn hồi căng ngang với đầu A cố định đang có sóng dừng. B là phần tử dây tại điểm bụng thứ hai tính từ đầu A, C là phần tử dây nằm giữa A và B. Biết A cách vị trí cân bằng của B và vị trí cân bằng của C những khoảng lần lượt là 30 cm và 5 cm, tốc độ truyền sóng trên dây là 50 cm/s. Trong quá trình dao động điều hoà, khoảng thời gian ngắn nhất giữa hai lần li độ của B có giá trị bằng biên độ dao động của C là
\choice
{$\dfrac{1}{15}\ s$}
{$\dfrac{2}{5}\ s$}
{$\dfrac{2}{15}\ s$}
{\True $\dfrac{1}{5}\ s$}
\loigiai{
\begin{itemize}
\item $AB=\dfrac{3\lambda}{4}=30\Rightarrow \lambda =40\ cm$.
\item C cách A 5\ cm $\Rightarrow AC=\dfrac{\lambda}{8}$.
\item Biên độ của C là: $A_C=2a\left|{cos\left(\dfrac{2\pi d}{\lambda}+\dfrac{\pi}{2}\right)}\right|=a\sqrt{2}$.
\item Khoảng thời gian ngắn nhất giữa hai lần li độ của B có giá trị bằng biên độ của C là: $\Delta t_{min}=\dfrac{T}4$.
\item Mặt khác: $v=50\ cm/s;\ \lambda =40\ cm\Rightarrow T=0,8\ s\Rightarrow \Delta t_{min}=\dfrac{T}4=\dfrac{1}5\ s$.
\end{itemize}}
\end{vidu} \haidong{20}
 